\chapter{激励轨迹 信号处理与摩擦补偿}\label{ch:excitation_friction}
\section{实验轨迹决定哪些参数能被看见}
辨识不是先随便录一段运动，再期待求解器把所有参数找出来。一个参数对输出的影响必须在数据中足够明显，并且能够与其他参数的影响区分，才有可能估计准确。这个要求叫激励充分性。反复做同一段慢速轨迹，可以降低部分随机噪声，却未必改善惯量与摩擦之间的区分。

考虑 $\tau=J_r\ddot q+F_v\dot q$。如果全部数据都满足 $\ddot q=2\dot q$，则 $\tau=(2J_r+F_v)\dot q$，只能识别组合 $2J_r+F_v$。让速度变化更复杂，使加速度与速度不总成比例，才能把这两项分开。多关节机械臂还需通过姿态变化区分重力、耦合惯量和不同连杆的贡献。

\section{有限项傅里叶激励的公式与导数}
\subsection{阶数就是保留多少个谐波}
对第 $i$ 个关节，采用
\begin{equation}\label{eq:fourier_di}
 q_i(t)=q_{i0}+\sum_{k=1}^{L}\left[
 \frac{a_{ik}}{k\omega_0}\sin(k\omega_0t)-\frac{b_{ik}}{k\omega_0}\cos(k\omega_0t)\right].
\end{equation}
$\omega_0=2\pi/T$ 是基角频率，$T$ 是一个周期，$L$ 是最高谐波序号；第 $k$ 项频率为 $k/T$。$a_{ik},b_{ik}$ 在这个参数化下的单位是角速度，不是角度。它与直接把角度幅值放在正余弦前面的写法等价，但系数单位不同。

逐项求导，分母中的 $k\omega_0$ 恰好抵消一次：
\begin{align}
 \dot q_i(t)&=\sum_{k=1}^{L}[a_{ik}\cos(k\omega_0t)+b_{ik}\sin(k\omega_0t)],\\
 \ddot q_i(t)&=\sum_{k=1}^{L}k\omega_0[-a_{ik}\sin(k\omega_0t)+b_{ik}\cos(k\omega_0t)],\\
 q_i^{(3)}(t)&=-\sum_{k=1}^{L}(k\omega_0)^2[a_{ik}\cos(k\omega_0t)+b_{ik}\sin(k\omega_0t)].
\end{align}
第三阶导数叫加加速度或 jerk。虽然只增加一个较小高频谐波，速度变化可能不大，但加速度与 jerk 因频率乘子而明显增大。

七阶与九阶在这个定义下分别保留 $7$ 与 $9$ 个谐波，最高频率分别为 $7/T$ 与 $9/T$；不是把时间 $t$ 提到七次或九次。“五次傅里叶级数”通常是口语不严格的说法，应回到公式确认其含义。若七个关节各用 $L$ 个谐波，未施加约束前有 $2\times7L$ 个正余弦系数，另有可能需要优化的中心位置。

\subsection{周期连续不等于起止静止}
取 $t=0$ 或 $t=T$，有
\[
 q_i(0)=q_{i0}-\sum_k\frac{b_{ik}}{k\omega_0},\qquad
 \dot q_i(0)=\sum_ka_{ik},\qquad
 \ddot q_i(0)=\omega_0\sum_k kb_{ik}.
\]
整数谐波保证这些量在一个周期后重复，但它们不自动为零。若需要从静止开始，应加上
\[
 \sum_ka_{ik}=0,\qquad\sum_kkb_{ik}=0.
\]
若还要求指定初始位置，必须同时调整 $q_{i0}$。有些实验会先用平滑过渡进入稳定周期，然后只采集中间完整周期；这也是合理设计，但过渡段不能当作满足周期条件的数据直接平均。

\begin{examplebox}[算例 9A 两个谐波形成起止静止的轨迹]
令 $T=2\ \mathrm s$，所以 $\omega_0=\pi\ \mathrm{rad/s}$；取 $a_1=0.2$、$a_2=-0.2$、$b_1=b_2=0$、$q_0=0$。有
\begin{align*}
 q(t)&=\frac{0.2}{\pi}\sin\pi t-\frac{0.1}{\pi}\sin2\pi t,\\
 \dot q(t)&=0.2\cos\pi t-0.2\cos2\pi t,\\
 \ddot q(t)&=-0.2\pi\sin\pi t+0.4\pi\sin2\pi t.
\end{align*}
初始速度 $0.2-0.2=0$，初始加速度为零，两个周期端点的位置也为零。$t=0.5\ \mathrm s$ 时，$q=0.0636620\ \mathrm{rad}$，$\dot q=0.2\ \mathrm{rad/s}$，$\ddot q=-0.628319\ \mathrm{rad/s^2}$。整个周期的位置峰值约 $0.0826993\ \mathrm{rad}$，速度绝对峰值 $0.4\ \mathrm{rad/s}$，加速度绝对峰值约 $1.71896\ \mathrm{rad/s^2}$。
\end{examplebox}

\begin{figure}[htbp]\centering
\includegraphics[width=0.89\linewidth]{figures/di_fourier_trajectory.pdf}
\caption{算例 9A 的位置、速度、加速度与 jerk。起止位置、速度、加速度连续，但这些边界条件本身不等于与任意外部静止段的 jerk 连续。}
\end{figure}

\section{轨迹优化到底在迭代什么}
优化变量是傅里叶系数。给定一组系数，就能算轨迹、回归矩阵和约束；优化器据此改变系数，再重新计算。所谓“迭代 164 次”，通常表示优化器执行了 164 次更新或主要迭代，具体计数还取决于软件；它不是机械臂已经来回运动 164 次，也不是辨识模型一定更新了 164 次。

设缩放白化后的矩阵为 $H_s$，信息矩阵为 $F=H_s^{\mathsf T}H_s$。常见目标有：
\begin{align}
 \text{条件数目标：}&\quad\min\ \frac{\sigma_{\max}(H_s)}{\sigma_{\min}(H_s)},\\
 \text{D 最优：}&\quad\max\ \log\det F,\\
 \text{A 最优：}&\quad\min\ \operatorname{tr}(F^{-1}),\\
 \text{E 最优：}&\quad\max\ \lambda_{\min}(F).
\end{align}
D 最优压小参数置信椭球的体积，A 最优压小方差总和，E 最优强化最弱参数方向。它们都依赖固定的噪声和参数尺度。D 最优在固定可逆参数坐标变换下的最优轨迹不变，因为行列式只增加一个与轨迹无关的常数因子；若每次轨迹都重新选缩放，则不能据此声称不受影响。

约束至少应对应实际模型中的位置、速度、加速度、jerk 或驱动器可实现性、力矩、碰撞和工作空间。离散网格上的约束满足不保证网格之间也满足；尤其高谐波轨迹，应使用更密的复核网格或连续时间上界。对碰撞距离，仅检查关节角极值更远远不够。

把轨迹整体放慢 $s$ 倍，即 $q_{\mathrm{slow}}(t)=q(t/s)$，则速度变为原来的 $1/s$，加速度变为 $1/s^2$，jerk 变为 $1/s^3$。这可以降低动态要求，却同时减弱惯量参数的激励；慢速试运行可以核对运动，但不能据此断定正式辨识数据仍有足够信息量。

\section{采样 时间同步与微分噪声}
\subsection{时间错位会制造不存在的模型误差}
如果力矩与状态相差 $\delta t$，一阶近似给出 $\tau(t+\delta t)-\tau(t)\approx\dot\tau(t)\delta t$。仅看重力部分，$\dot g=(\partial g/\partial q)\dot q$。若某工作点的重力梯度 $40\ \mathrm{N\,m/rad}$、速度 $1\ \mathrm{rad/s}$、时间错位 $2\ \mathrm{ms}$，伪力矩误差就是 $40\times1\times0.002=0.08\ \mathrm{N\,m}$。

时间戳应明确表示采样时刻、发送时刻还是接收时刻。互相关可帮助发现大致错位，但峰值不为零也可能来自执行器动态或模型误差；不能不分析物理链条就通过任意平移数据把残差调到最小。

\subsection{求导为什么放大量化误差}
设每个位置样本独立带有均匀量化误差，量化间隔为 $\Delta$，则位置误差方差为 $\Delta^2/12$。一阶后向差分
\[
 \hat v_k=\frac{q_k-q_{k-1}}{T_s}
\]
使两个独立误差相减，方差相加，因此
\[
 \sigma_v^2=\frac{2(\Delta^2/12)}{T_s^2},\qquad
 \sigma_v=\frac{\Delta}{\sqrt6T_s}.
\]
二阶差分 $(q_{k+1}-2q_k+q_{k-1})/T_s^2$ 的三个噪声系数平方和为 $1+4+1=6$，所以
\[
 \sigma_a=\frac{\Delta}{\sqrt2T_s^2}.
\]
加速度噪声按采样周期的平方反比增长。更快采样提供更多信息，但直接使用未经处理的二阶差分，并不自动得到更好的加速度。

\begin{examplebox}[算例 9B 十七位编码器的差分误差量级]
一圈 $2^{17}$ 个计数时，$\Delta=2\pi/2^{17}=4.79369\times10^{-5}\ \mathrm{rad}$。若 $T_s=1\ \mathrm{ms}$，则 $\sigma_v\approx0.01957\ \mathrm{rad/s}$，$\sigma_a\approx33.90\ \mathrm{rad/s^2}$。若 $T_s=5\ \mathrm{ms}$，则分别约为 $0.003914\ \mathrm{rad/s}$ 和 $1.356\ \mathrm{rad/s^2}$。这只是独立量化噪声模型的估计，真实误差还可能有相关性、机械振动和编码器周期误差。
\end{examplebox}

\subsection{可以怎样估计速度和加速度}
离线可对位置做平滑后差分、Savitzky--Golay 局部多项式拟合，或在周期实验中拟合谐波再解析求导。局部多项式的本质是在一个窗口内拟合 $q(t)\approx c_0+c_1t+c_2t^2+\cdots$，窗口中心速度为 $c_1$、加速度为 $2c_2$。窗口越宽通常越降噪，却越难保留快速真实变化。

谐波拟合则使用已知频率的正余弦基函数，估计实际测量位置的系数。应拟合测得的位置，不能直接把指令轨迹当成真实运动；否则跟踪滞后和驱动器饱和会被错误归到模型参数。真实运动包含未建模谐波、漂移或非周期过渡时，过低的谐波阶数也会产生偏差。

\section{为什么力矩和回归矩阵必须同滤波}
假设 $y(t)=H(t)\beta+\epsilon(t)$，参数 $\beta$ 在这段实验中恒定。对任意线性滤波算子 $\mathcal F$，有
\begin{equation}
 \mathcal F[y]=\mathcal F[H\beta]+\mathcal F[\epsilon]
 =(\mathcal F[H])\beta+\mathcal F[\epsilon].
\end{equation}
因此可在力矩和每一列最终回归特征上应用相同滤波。这里包括惯性列、重力列以及摩擦列，后续使用同一套边缘样本掩码。

\begin{pitfallbox}[先滤位置再造矩阵 不等于对回归矩阵滤波]
一般有 $H(\mathcal F[q])\ne\mathcal F[H(q)]$，因为 $H$ 含正弦、余弦、速度乘积和符号函数等非线性。位置平滑用于得到可靠状态；构造最终回归矩阵后，对其各列和力矩同滤波，是另一件事。只把位置和力矩使用同一个滤波器，并不足以自动保证最终辨识方程两边一致。
\end{pitfallbox}

周期运动可能在回归矩阵里产生更高频率。例如 $\dot q=\sin\omega t$，则 $\dot q^2=(1-\cos2\omega t)/2$，多出二倍频；$\cos(q(t))$ 更可能包含多个谐波。故滤波截止频率不能只根据位置轨迹最高谐波机械确定。

\subsection{四阶双向低通到底是什么意思}
以四阶 Butterworth 数字低通为例，单向传递响应为 $H_f(e^{j\Omega})$。前向再反向滤波，在远离边界且近似平稳的内部区间，其有效频率响应是
\[
 H_{\mathrm{fb}}(e^{j\Omega})=H_f(e^{j\Omega})H_f(e^{-j\Omega})=|H_f(e^{j\Omega})|^2.
\]
相位相互抵消，但幅值也被再乘一次。单向设计截止频率处幅值为 $1/\sqrt2$，双向后为 $1/2$，即约 $-6.02\ \mathrm{dB}$，不是 $-3.01\ \mathrm{dB}$。说“等效八阶”仅指渐近幅度滚降，不意味着它是一个可实时运行的普通因果八阶滤波器。

\begin{examplebox}[算例 9C 二百赫兹采样下八赫兹双向滤波]
取 $f_s=200\ \mathrm{Hz}$，用四阶数字 Butterworth 设计单向截止频率 $8\ \mathrm{Hz}$。计算双向有效幅值，在 $4,8,12,16\ \mathrm{Hz}$ 处分别约为
\[
 0.99623,\quad0.50000,\quad0.035674,\quad0.0034228.
\]
因此 $8\ \mathrm{Hz}$ 的真实信号已经减半，$12\ \mathrm{Hz}$ 成分只剩约 $3.57\%$。这能显著降噪，也可能删除与惯量或柔性有关的信息。该数字说明的是一组指定设计，不是所有“八赫兹低通”的共同表现。
\end{examplebox}

\begin{figure}[htbp]\centering
\includegraphics[width=0.92\linewidth]{figures/di_filter_eight_hz.pdf}
\caption{指定采样率和设计方法下的四阶单向与双向滤波响应。双向滤波无内部相位延迟，但幅值平方、边缘效应和非因果性仍需处理。}
\end{figure}

选 $8\ \mathrm{Hz}$ 是否合适，应比较几个候选截止频率下的真实信号频谱、回归列能量、奇异值、参数稳定性和独立验证误差。若参数随截止频率大幅变化，应怀疑模型、噪声、对齐或高频真实动态。双向滤波需要未来样本，仅用于离线处理；在线控制应使用因果滤波并把其相位延迟计入闭环。

\section{不显式求加速度的积分回归}
前面的方法通过平滑得到加速度。另一条思路是先把动力学在时间窗口内积分，把加速度项变成速度差。对已扣除重力的单关节模型
\[
 \tau=J_r\dot v+F_vv+F_c\operatorname{sgn}(v),
\]
从 $t_0$ 积分到 $t_1$，有
\begin{equation}
 \int_{t_0}^{t_1}\tau\,dt
 =J_r[v(t_1)-v(t_0)]+F_v[q(t_1)-q(t_0)]
 +F_c\int_{t_0}^{t_1}\operatorname{sgn}(v)\,dt.
\end{equation}
未知参数依然线性，但回归量不再包含显式加速度。若窗口内始终正向运动，最后一个积分就是窗口时长；若发生换向，则必须根据速度符号分别计算。重力未扣除时，还需保留 $h g_0\int\cos q\,dt$，不能因为用了积分模型就把重力删掉。

\begin{examplebox}[算例 9G 用一个时间窗口核对积分动力学]
取 $J_r=0.1$、$F_v=0.2$、$F_c=0.3$，窗口时长 $0.2\ \mathrm s$，速度以常加速度由 $0.2$ 增至 $0.4\ \mathrm{rad/s}$。速度差为 $0.2$，位置增量为平均速度乘时间 $0.3\times0.2=0.06\ \mathrm{rad}$，全程正向。因此
\[
 \int\tau dt=0.1\times0.2+0.2\times0.06+0.3\times0.2=0.092\ \mathrm{N\,m\,s}.
\]
直接按瞬时力矩计算：加速度为一，力矩从 $0.1+0.2\times0.2+0.3=0.44$ 线性增至 $0.48$，平均值 $0.46$，乘以窗口长度同样得到 $0.092$。
\end{examplebox}

一个窗口只提供一条方程，不能单独估出三个参数。不同窗口还需具有独立的速度变化、位移和持续时间特征。相邻重叠窗口的噪声彼此相关，窗口太长也会把不同参数影响混在一起，并累积传感器偏置。

多关节系统不能简单把 $M(q)\ddot q$ 积成 $M\Delta\dot q$，因为 $M$ 随时间改变。正确分部积分含有
\[
 \int M\ddot q\,dt=[M\dot q]_{t_0}^{t_1}-\int\dot M\dot q\,dt.
\]
必须把额外项与原有速度项一致处理，或从已知线性回归结构推导滤波动力学、动量或能量回归。积分法降低了对直接加速度的依赖，却没有免除对模型、时间同步和参数可辨识性的要求。

\section{库仑与黏性摩擦怎样从曲线中读出来}
对非零速度，最小模型为
\begin{equation}
 \tau_f(v)=F_c\operatorname{sgn}(v)+F_vv.
\end{equation}
在正速度一侧，它是一条截距 $F_c$、斜率 $F_v$ 的直线；负速度一侧截距为 $-F_c$。库仑摩擦近似与速度大小无关，黏性摩擦随速度线性增加。这里 $\tau_f$ 放在动力学左侧，耗散要求 $\tau_fv\ge0$。

\begin{examplebox}[算例 9D 通过三种速度辨识摩擦直线]
在正速度 $0.2,0.5,1.0\ \mathrm{rad/s}$ 下，扣除其他动力学项后的平均阻力分别为 $0.32,0.35,0.40\ \mathrm{N\,m}$。任取第一与第三点，斜率
\[
 F_v=\frac{0.40-0.32}{1.0-0.2}=0.10\ \mathrm{N\,m\,s/rad}.
\]
截距 $F_c=0.32-0.10\times0.2=0.30\ \mathrm{N\,m}$。第二点预测为 $0.30+0.10\times0.5=0.35$，恰好吻合。有噪声时应对所有点一起做最小二乘，并给出重复测量离散程度。
\end{examplebox}

\subsection{正反向半和半差为什么有条件}
同一姿态、相同速度绝对值、加速度近似为零，且摩擦对称时：
\[
 \tau_+=g+F_c+F_vv,\qquad\tau_-=g-F_c-F_vv.
\]
因此 $g=(\tau_++\tau_-)/2$，$F_c+F_vv=(\tau_+-\tau_-)/2$。若 $g=2.5$，$F_c=0.3$，$F_v=0.1$，$v=0.5$，则 $\tau_+=2.85$、$\tau_-=2.15$，半和为 $2.5$、半差为 $0.35\ \mathrm{N\,m}$。

如果正反向数据来自不同角度，重力项可能不同；如果电缆力矩和摩擦不对称，半和还包含其偶对称部分；如果其他关节在运动，速度耦合也可能未消失。单姿态静止电流更不能直接视为准确重力，因为静摩擦可以承担一部分保持力矩。

\section{静摩擦不能简单等同于符号函数在零点取零}
静止时，摩擦不是固定输出一个数，而是在允许范围内恰好抵消外加的剩余驱动力。简化写作
\begin{equation}
 v=0:\quad\tau_f\in[-F_s,F_s],\qquad
 \tau_f=\tau_{\mathrm{cmd}}+\tau_{\mathrm{ext}}-g\quad\text{若静止能保持}.
\end{equation}
$F_s$ 是起动阈值。只有当所需保持摩擦超过这个区间，物体才开始滑动。把 $\operatorname{sgn}(0)$ 设为零，是某个数值函数的定义，不是静摩擦物理规律。

例如 $F_s=0.5\ \mathrm{N\,m}$、扣除重力后的驱动力为 $0.3\ \mathrm{N\,m}$，关节仍可能不动；实际静摩擦此时是 $0.3$，不是 $0.5$，也不是零。这直接解释低刚度控制中的回正死区：恢复力矩变小到静摩擦能平衡的范围，系统可能停在非零误差处。

Stribeck 曲线描述从起动到滑动过程中，摩擦随低速增加先下降的现象。一种常见滑动模型为
\begin{equation}
 \tau_f(v)=\left[F_c+(F_s-F_c)e^{-(|v|/v_s)^\alpha}\right]\operatorname{sgn}(v)+F_vv,\qquad v\ne0.
\end{equation}
$v_s$ 是过渡速度尺度，$\alpha$ 控制曲线形状。这个静态速度函数仍不能独立描述真正的预滑动记忆、迟滞和粘滑循环。

\begin{examplebox}[算例 9E 低速处摩擦并不一定更小]
取 $F_s=0.5$、$F_c=0.3\ \mathrm{N\,m}$、$F_v=0.1\ \mathrm{N\,m\,s/rad}$、$v_s=0.05\ \mathrm{rad/s}$、$\alpha=2$。在 $v=0.02$ 时
\[
 \tau_f=0.3+0.2e^{-0.16}+0.002=0.472429\ \mathrm{N\,m}.
\]
在 $v=0.1$ 时
\[
 \tau_f=0.3+0.2e^{-4}+0.01=0.313663\ \mathrm{N\,m}.
\]
速度增加五倍，摩擦反而先降低。这一段负斜率与弹性、惯性共同作用，可能形成走停交替的粘滑现象。
\end{examplebox}

\begin{figure}[htbp]\centering
\includegraphics[width=0.91\linewidth]{figures/di_friction_models.pdf}
\caption{库仑加黏性、Stribeck 与平滑补偿曲线。零速竖直线段表示静摩擦允许的集合，而非零速时只有一个固定摩擦值。}
\end{figure}

若必须描述预滑动与记忆，可引入 LuGre 等动态摩擦模型。一个常见形式是
\begin{align}
 \dot z&=v-\frac{|v|}{g_z(v)}z,\\
 \tau_f&=\sigma_0z+\sigma_1\dot z+\sigma_2v.
\end{align}
$z$ 表示微观接触的平均弹性变形，$g_z(v)$ 是由稳态摩擦曲线确定的长度尺度，$\sigma_0$ 为微接触刚度，$\sigma_1$ 为微阻尼。若稳态摩擦幅值为 $F(v)$，可令 $g_z(v)=F(v)/\sigma_0$。必须说明这一单位关系，否则不同文献中 $g(v)$ 是力还是长度会混淆。动态模型新增状态和参数，应由数据支持其必要性，不能仅因名称更复杂就直接使用。

\section{联合估计与交替估计怎样配合}
设模型 $y=Y\pi+F\gamma$，其中 $\pi$ 为惯性参数，$\gamma$ 为线性摩擦参数。联合最小二乘直接解
\[
 \min_{\pi,\gamma}\|W(y-[Y\ F][\pi;\gamma])\|^2.
\]
如果先忽略摩擦拟合惯性，摩擦残差可能投影到 $Y$ 的列空间；之后用残差估摩擦，再固定摩擦重估惯性，可以逐步减轻相互污染。但不能保证一两轮就收敛，也不能通过迭代自动创造缺失的信息。

对一个列满秩的纯线性凸二次问题，精确的交替块最小化可收敛到联合最优解；如果惯性与摩擦特征强相关，可能收敛很慢。若包含 $v_s$ 等非线性参数、数据剔除规则变化或饱和模型，甚至可能出现局部极小或不稳定的反复更新。因此，线性部分优先联合解，低速非线性部分再利用专门实验估计，并在同一验证集比较。

\section{补偿不是把拟合曲线原封不动全部加回去}
为了减少速度噪声引发的换向跳变，可用
\[
 \tau_{\mathrm{ff},f}=\gamma\left[F_c\tanh(v/\varepsilon)+F_vv\right],\qquad0\le\gamma\le1.
\]
$\varepsilon$ 是速度平滑尺度，$\gamma$ 是补偿比例。这能让零附近力矩连续，但不会准确再现静摩擦集合。$\varepsilon$ 越小，零附近斜率 $\gamma F_c/\varepsilon$ 越大，测量噪声更容易引发明显补偿变化。

若补偿基于实测速度，且模型估计过大，它会沿运动方向主动注入过多能量，表现得像负阻尼。若补偿基于期望速度，停滞时又可能与真实运动方向不一致。具体选择要结合控制结构、速度估计和目标工况；补偿后应观察力矩、速度、能量与温度变化，而不只看一次回正误差。

\begin{examplebox}[算例 9F 平滑尺度决定零附近的敏感度]
取 $F_c=0.3\ \mathrm{N\,m}$、$\gamma=0.8$、$\varepsilon=0.02\ \mathrm{rad/s}$。零附近斜率为 $0.8\times0.3/0.02=12\ \mathrm{N\,m\,s/rad}$。若把 $\varepsilon$ 缩小到 $0.002$，斜率增至 $120$。同样的 $0.001\ \mathrm{rad/s}$ 速度噪声，在局部线性近似下可引起约 $0.012$ 与 $0.12\ \mathrm{N\,m}$ 的补偿变化，差十倍。
\end{examplebox}

\section{面向七轴低阻抗的结果应该怎样记录}
辨识交付物至少应明确：几何和负载版本、关节顺序与符号、反馈力矩定义、时间对齐、温度范围、滤波器完整设计、样本剔除条件、基参数变换、噪声白化与列缩放、训练及验证轨迹、各关节绝对和相对误差。条件数应附矩阵定义，不能只留下一个“40”。

腕部关节力矩小，电流偏置、线缆拖拽和负载质心误差可能占更大比例。末端误差还要经过雅可比与刚度映射，不能仅凭“J5--J7 摩擦估计大”就认定机械摩擦异常。应把残差按速度、姿态、加速度、温度和运动方向分别作图，确认其结构，再决定增加哪一类模型。

\section{计算练习与分步解答}
\subsection*{练习 9.1 七阶与九阶的最高频率}
周期 $T=10\ \mathrm s$，分别使用七阶、九阶傅里叶级数。求基频和最高频率。

\textbf{解答。}基频 $f_0=1/T=0.1\ \mathrm{Hz}$。七阶最高 $7f_0=0.7\ \mathrm{Hz}$，九阶最高 $0.9\ \mathrm{Hz}$。如果把周期改成 $2\ \mathrm s$，两者变为 $3.5$ 和 $4.5\ \mathrm{Hz}$。阶数本身不能单独说明轨迹快慢。

\subsection*{练习 9.2 同一轨迹放慢两倍}
某轨迹速度、加速度、jerk 峰值为 $0.8$、$3$、$20$，单位依次为 $\mathrm{rad/s}$、$\mathrm{rad/s^2}$、$\mathrm{rad/s^3}$。整体放慢两倍后是多少？

\textbf{解答。}时间缩放 $s=2$，三者分别除以 $2$、$4$、$8$，得到 $0.4$、$0.75$、$2.5$。位置范围不变，动态惯性力矩却通常显著下降。

\subsection*{练习 9.3 从正反向数据分离对称摩擦}
同一姿态、速度 $\pm0.4\ \mathrm{rad/s}$ 时，稳态力矩分别为 $3.7$ 和 $2.9\ \mathrm{N\,m}$。在对称模型假设下求重力与摩擦幅值。

\textbf{解答。}半和 $(3.7+2.9)/2=3.3$，半差 $(3.7-2.9)/2=0.4\ \mathrm{N\,m}$。这次数据只确定 $F_c+0.4F_v=0.4$，不能分别确定两者；还需要其他速度。

\subsection*{练习 9.4 低刚度下的静摩擦误差范围}
单关节恢复力矩为 $-K_pe$，静摩擦阈值 $F_s=0.4\ \mathrm{N\,m}$，$K_p=8\ \mathrm{N\,m/rad}$，无其他偏置。求可能静止的误差区间。

\textbf{解答。}只要 $|K_pe|\le F_s$，静摩擦就可能平衡恢复力矩。所以 $|e|\le0.4/8=0.05\ \mathrm{rad}=2.865^\circ$。这是一个理想静态可能区间，不是预测每次一定停在其边缘。若存在重力模型残差，区间中心还会移动。

\section*{本章与原稿的对应}
本章整合 [A，第33--34页] 与 [B，第42--53页]，重点补足傅里叶边界条件、优化迭代的含义、差分噪声推导、同滤波的线性代数依据、双向滤波的有效幅值以及摩擦分离的假设。原稿中的固定截止频率倍数、固定条件数目标和固定收敛轮数均改为需要由模型和数据判断的设计选择。
